\section{第二部分 系统组成及功能说明（一）：需求分析与设计约束}

\subsection{材料实验痛点}
近红外荧光粉等材料体系通常存在候选配方多、物理机制复杂、真实烧制成本高的问题。研究员往往需要先提出一批候选，再逐个经历称量、研磨、烧制、XRD、PL 光谱和结果判断。若每轮要尝试 15-20 个候选，人工决策和实验等待会成为主要瓶颈。

本作品把需求拆成三个层次：第一层是 AI 预测，目的是在真实实验前快速排除高风险候选；第二层是具身执行，目的是让机器人承担可规范化的取放和转移动作；第三层是数据回流，目的是把实验结果变成下一轮预测的证据，而不是散落在截图、表格和人工记忆里。

\subsection{功能需求}
\begin{table}[H]
\centering
\caption{系统功能需求与对应实现。}
\begin{tabular}{p{3.5cm}p{5.4cm}p{5.8cm}}
\toprule
需求 & 设计实现 & 初赛状态 \\
\midrule
配方预筛 & AI 脑解析配方，调用物理模型、实测类比、本地模型和规则库 & 已跑通 \\
XRD/PL 分析 & 四条表征线分别处理视觉图像和数值谱图 & 已跑通 \\
移动感知 & 车载脑融合 LiDAR、深度、里程计并构建 SLAM 地图 & 已跑通 \\
影子规划 & Lab-FSD 输出 BEV、future BEV、policy token、risk 和 safety gate & 在线 shadow \\
末端执行 & F407 控制推杆、舵机、电磁铁和短行程升降动作 & 安全降级 \\
固定工位取放 & arm01 末端相机识别工位，AI 脑确认后固定点取袋 & 已实测 \\
证据展示 & 公网站展示状态、证据、API、发布和演示剧场 & 已部署 \\
\bottomrule
\end{tabular}
\end{table}

\subsection{嵌入式约束}
项目选择 RDK X5 的原因不是单纯追求模型参数规模，而是要在边缘设备上同时管理摄像头、BPU、CPU、本地服务、ROS2、传感器、机械执行和网络证据。设计中必须处理以下约束：
\begin{itemize}
  \item X5 不适合在线运行重型材料模拟或训练任务，MACE、MatterGen、LoRA 训练等放在 PC/云端预处理，X5 读取缓存和部署产物。
  \item BPU 的算子、CMA 和模型切段限制决定了 Transformer 只能以 slot / probe 方式工程展示，不能夸大为长文本实时生成。
  \item 移动底盘和升降台属于安全敏感执行器，初赛阶段必须采用安全员接管和硬件降级策略。
  \item 公网站点只能做状态展示和证据说明，不能下发底盘、机械臂、升降台、电磁铁等危险动作。
\end{itemize}

\begin{warnbox}
\textbf{安全原则：} 本作品报告中的“自主”“规划”“世界模型”均指算法在线运行和诊断输出。未完成全量安全验证的执行链路不写成无人值守闭环。
\end{warnbox}
